\paragraph{Integral averaging kernels and the projected lattice.}
\teamanchor{proof.integral-ramified-kernel}%
The complex dimension calculation has the following integral
counterpart. It also distinguishes two lattices with the same
rational span. Let \(\mathcal O\) be the ring of integers of a number
field \(K\), let \(q\ge2\) and \(f\ge1\) be integers, and put
\[
 A_q=\ker(\Sigma:\mathcal O^q\longrightarrow\mathcal O),
 \qquad \Sigma(x)=\sum_{i=1}^q x_i.
\]
Suppose a definite class lattice has an integral decomposition
\[
 M=\bigoplus_{\omega\in\Omega}
       (\mathcal O^q)^{\otimes f}\otimes_{\mathcal O} L_\omega,
 \qquad r=\sum_\omega\operatorname{rk}L_\omega>0,
\]
where the \(L_\omega\) are finite projective modules. At every
complex embedding require the class-coordinate copies of each
\(L_\omega\) to be orthogonal with identical Hermitian metrics;
different orbits are orthogonal too. A common positive metric
factor on one orbit is allowed. The partial summation map \(s_j\)
sums over the \(j\)-th \(q\)-element coordinate, leaving all the
other coordinates unchanged. Define
\[
 H=\bigcap_{j=1}^f\ker s_j\subset M,
 \qquad d_M=q^f r,\qquad d_r=(q-1)^f r.
\]
Then \(H\) is saturated, has rank \(d_r\), commutes with arbitrary
extension of the coefficient ring, and satisfies
\begin{equation}
\label{eq:integral-ramified-kernel}
 H=\bigoplus_\omega A_q^{\otimes f}\otimes L_\omega,
 \qquad
 \mu(H)=\mu(M)-\frac{f\log q}{q-1},\qquad
 \mu_{\max}(H)\le\mu_{\max}(M).
\end{equation}
The slopes use the squared metrics of the surrounding argument.

To prove the integral assertion, use
\(\mathcal O^q=A_q\oplus\mathcal O e_q\). Expanding the tensor
product by this direct sum, the kernel of \(s_j\) consists exactly
of the summands with \(A_q\) in position \(j\). Their intersection
has \(A_q\) in every position. All the complementary summands are
projective, proving saturation and the assertion after any base
change. No division by \(q\) enters this argument.

In one factor the basis \(a_i=e_i-e_q\), \(1\le i<q\), of
\(A_q\) has Gram matrix
\(G=I+\mathbf1\mathbf1^{\mathsf T}\), with determinant \(q\).
Writing \(a=q-1\), the identity
\[
 \det(G^{\otimes f})=q^{f a^{f-1}}
\]
shows that the additional log squared covolume per retained rank
is \(f\log q/a\). The contribution per rank from the fibre
metric and its finite determinant is the same before and after
restriction. Each orbit retains the same rank ratio \((a/q)^f\),
so their weighted means retain the same weights. This proves
the mean formula. Every rational subspace of \(H\otimes K\)
has the same induced lattice and metric as its intersection with
\(M\), since \(H\) is saturated. Its slope is bounded by
\(\mu_{\max}(M)\), proving the last inequality.

\paragraph{The denominator introduced by orthogonal projection.}
\teamanchor{proof.ramified-projector-index}%
Let \(e_j=I-q^{-1}J\) on the \(j\)-th factor, where every entry
of \(J\) is one, and let \(e=\prod_j e_j\) on \(M\otimes K\).
This is the orthogonal projection onto \(H\otimes K\) in the
specified metrics. The projected lattice \(P=eM\) satisfies
\begin{equation}
\label{eq:ramified-projector-index}
 H\subseteq P\subseteq q^{-f}H,\qquad
 \operatorname{length}_{\mathcal O_v}(P_v/H_v)
     =\frac{f d_r}{q-1}\operatorname{ord}_v(q)
\end{equation}
at every finite place \(v\), with \(\operatorname{ord}_v\) the
integer-valued DVR valuation. In particular \(P_v=H_v\) when
\(q\) is a unit.

Indeed, \(b_i=e_i-q^{-1}\mathbf1\), \(1\le i<q\), form a basis
of the one-factor projected lattice: the last projected coordinate
is \(-\sum_i b_i\). The relation
\(a_j=b_j+\sum_i b_i\) gives inclusion matrix \(G\) from the
kernel lattice to the projected lattice. It has determinant \(q\).
Tensoring the inclusions and choosing local bases of the
\(L_\omega\) gives determinant valuation
\(r f(q-1)^{f-1}\operatorname{ord}_v(q)\), proving the length.
The first inclusion follows also from \(e|_H=I\). Since
\(q^f eM\) is integral and killed by every \(s_j\), it lies in
\(H\), proving the other inclusion. The Gram matrix of the
\(b_i\) is \(I-q^{-1}\mathbf1\mathbf1^{\mathsf T}=G^{-1}\).
Consequently the projected lattice has the opposite mean correction:
\[
 \mu(P)=\mu(M)+\frac{f\log q}{q-1}.
\]
This calculation concerns the finite class coordinates, and is a
direct lattice calculation rather than an additional candidate
for new mathematics.

For the word-prime model take \(q=T_w-1\). The possible difference
between these two lattices is supported on primes dividing
\(T_w-1\). At places over \(T_w\) the two model lattices coincide.
This does not determine the index of the actual source \(N_w\)
there. Nor does equality of the complex spans identify \(N_w\)
with either model. The required integral orbit decomposition and
its metrics must themselves be identified with those of the
actual definite class lattice.

\paragraph{A local test for the remaining donor comparison.}
\teamanchor{proof.ramified-donor-matrix-test}%
Fix a single \(K\)-linear identification
\(\phi:N_w\otimes K\longrightarrow H_w\otimes K\) and set
\(q=T_w-1\) in the preceding model. At each finite place choose
bases and let \(C_v\in\mathrm{GL}_{d_r}(K_v)\) be the matrix
whose columns generate \(\phi((N_w)_v)\) in an \((H_w)_v\)-basis.
Let \(\alpha_{1v}\le\cdots\le\alpha_{d_rv}\) be its elementary
divisor valuations; these integers can be negative before
integrality is proved. Then the two conditions away from \(T_w\)
in the donor comparison, with \(e_v\) the absolute ramification
index and \(\tau_v\) a nonnegative integer, are exactly
\begin{equation}
\label{eq:ramified-donor-matrix-test}
 \alpha_{1v}\ge e_v\tau_v,\qquad
 \sum_{i=1}^{d_r}\alpha_{iv}
       \le e_v d_M(\tau_v+\iota_v).
\end{equation}
At \(v\mid T_w\) the corresponding conditions are
\(\alpha_{1v}\ge0\) and the stated weighted level-index bound
with \(e_v\ell_{wv}=\sum_i\alpha_{iv}\). Indeed Smith normal
form gives the containment from the least exponent and, once
all exponents are nonnegative, the quotient length from their
sum. This proves both directions of the test.

The same \(\phi\) must be an isometry at every complex embedding.
In compatible bases this is the identity
\(G_{N,\sigma}=C_\sigma^*G_{H,\sigma}C_\sigma\), with
\(C_\sigma\) the matrix of that same map, not an independently
chosen local identification. The matrix test makes the residual
claim explicit; matrices satisfying it for the actual selected
sources and their metrics have not been supplied here.
